%% A simple template for a lab or course report using the Hagenberg setup
%% based on the standard LaTeX 'report' class

%%% Magic comments for setting the correct parameters in compatible IDEs
% !TeX encoding = utf8
% !TeX program = pdflatex 
% !TeX spellcheck = en_US
% !BIB program = biber

\RequirePackage[utf8]{inputenc} % Remove when using lualatex or xelatex!
\RequirePackage{hgbpdfa}				% Remove if NO PDF/A output is desired

\documentclass[language=english,titlepage=false,smartquotes=true]{hgbreport}
% Supported options in [..]:
%    Main language: language=english (default), language=german
%    Typographic quotes: smartquotes=true, smartquotes=false (default)
%    APA citation style: apa=true, apa=false (default)
%    Title page: titlepage=true (default), titlepage=false
%    Page layout: twoside=false (single-sided, default), twoside=true (double-sided)
%    Update check: updatecheck=true (default), updatecheck=false (to suppress)
%%%-----------------------------------------------------------------------------

\renewcommand{\chapter}[1]{} % Disable \chapter command
\graphicspath{{images/}}     % Location of images and graphics
\bibliography{hgbreferences}    % Biblatex bibliography file (hgbreferences.bib)
\ExecuteBibliographyOptions{backref=false} % No back references in bibliography

%%%-----------------------------------------------------------------------------
\setcounter{chapter}{1}	% <----- set to assignment number!
\setcounter{secnumdepth}{4}  % Nummerierung bis Ebene 4

%%%-----------------------------------------------------------------------------
\begin{document}
%%%-----------------------------------------------------------------------------

\author{Niklas Etzlinger and Rupert Parzmair}                      % Your name
\title{VCO2IL: Visual Computing (IM.ma VZ SS26)\\ % Name of the course
				Lab Report \arabic{chapter}}
\date{\today}

%%%-----------------------------------------------------------------------------
\maketitle
%%%-----------------------------------------------------------------------------

\begin{abstract}\noindent
ToDO: Write a brief abstract (max. 150 words) summarizing the main goals, methods and results of this assignment.
\end{abstract}

%%%-----------------------------------------------------------------------------

\section{Related Work: Object Detection in Aerial and Satellite Imagery}


\subsection{Overview}

Object detection in aerial and satellite imagery refers to the automatic localization and classification of objects of interest in overhead images of the Earth's surface. Typical target objects include vehicles, ships, airplanes, buildings, storage tanks, roads, and other man-made structures. This task is an important component of remote sensing image interpretation and supports a wide range of applications, such as urban planning, geographic information system updating, disaster response, precision agriculture, environmental monitoring, and security-related analysis \cite{li2020object,ding2021object}.

Although object detection has been widely studied in natural images, aerial and satellite imagery introduces several domain-specific characteristics that make the task fundamentally different from ground-level detection. Natural images usually capture objects from a horizontal perspective and show their profile information, whereas remote sensing images are acquired from an overhead perspective and mainly capture roof or top-view information \cite{li2020object}. As a result, object appearance, scale, orientation, and surrounding context differ substantially from standard computer vision datasets. In addition, aerial images often contain objects with arbitrary orientations, large scale variations, dense object distributions, complex backgrounds, and very large image sizes \cite{ding2021object}.

These properties make direct transfer of standard object detectors to aerial and satellite imagery difficult. For example, small vehicles may occupy only a few pixels, while airports or harbors can span large regions of an image. Objects such as ships, airplanes, and vehicles are often not aligned with the image axes, which motivates the use of oriented bounding boxes instead of conventional horizontal bounding boxes \cite{ding2021object}. Furthermore, large satellite or aerial images often need to be processed using tiling or sliding-window strategies because they cannot be passed through detection networks at full resolution.

This chapter reviews the main challenges and methods for object detection in aerial and satellite imagery. First, it discusses important challenges such as scale variation, arbitrary object orientation, dense object distributions, large image sizes, and foreground-background imbalance. Second, it introduces oriented object detection and explains why horizontal bounding boxes are often insufficient for overhead imagery. Finally, relevant benchmark datasets and evaluation metrics are summarized, with particular attention to widely used datasets such as DOTA, DIOR, and xView.

\subsection{Challenges in Aerial Imagery}

A major challenge in aerial and satellite imagery is the strong variation in object scale. Small objects such as cars or small vehicles may occupy only a few pixels, while large structures such as airports, harbors, bridges, or industrial buildings can cover large image regions. This makes single-scale detection insufficient, because a feature representation that is suitable for large objects often lacks the spatial detail needed for small objects. Feature Pyramid Networks (FPN) address this problem by constructing multi-scale feature representations inside the network, allowing detections to be made at different feature levels \cite{ding2021object,lam2018xview,lin2017feature}.

Another important difference compared to ground-level imagery is the arbitrary orientation of objects. In aerial images, objects such as ships, airplanes, vehicles, or storage tanks are observed from a top-down perspective and are therefore not necessarily aligned with the horizontal or vertical image axes. Conventional horizontal bounding boxes (HBBs) can include large amounts of background for rotated objects and may localize them only imprecisely. For this reason, oriented bounding boxes (OBBs) are often more appropriate in aerial object detection, since they can represent the actual object extent more accurately \cite{ding2021object,ding2018learning}.

Dense object distributions further complicate detection. In scenes such as parking lots, harbors, airports, or urban traffic areas, many objects can appear close to each other or even overlap in the image. In such cases, horizontal boxes may cover multiple neighboring instances, which can lead to ambiguous feature extraction and localization errors. Moreover, post-processing methods such as non-maximum suppression become more difficult in high-density scenes, because correct detections of nearby objects may have high overlap and can be suppressed incorrectly \cite{ding2021object,ding2018learning}.

Large image sizes are another practical issue in aerial and satellite imagery. Remote sensing images can be much larger than typical natural images, and in Earth observation settings image sizes may reach tens of thousands of pixels in width and height. Such images cannot usually be processed directly by standard convolutional object detectors because of memory and computation limits. Therefore, large images are commonly divided into smaller tiles or processed using sliding-window strategies. However, this introduces trade-offs: overlapping tiles reduce the risk of cutting objects at tile borders but increase computation time, while smaller tiles may lose important contextual information \cite{ding2021object}.

Finally, aerial imagery often suffers from a strong foreground-background imbalance. Most pixels in an overhead image usually belong to background regions such as roads, fields, water, rooftops, or bare ground, while the target objects occupy only a small fraction of the image. This imbalance is especially problematic for dense one-stage detectors, where many easy background examples can dominate the training process. Focal Loss addresses this issue by down-weighting well-classified examples and focusing the training more strongly on hard examples, which makes it a suitable theoretical solution for foreground-background imbalance in dense object detection \cite{lin2017focal}.


\subsection{Oriented Object Detection}

\subsubsection{Limitations of Horizontal Bounding Boxes}

Horizontal bounding boxes (HBBs) are the standard representation in general object detection, but they are often insufficient for aerial imagery. Since objects in overhead images can appear at arbitrary angles, an axis-aligned box may not match the actual object shape. For example, a rotated ship or vehicle can only be enclosed by a much larger horizontal rectangle, which includes a large amount of irrelevant background. This reduces the consistency between the visual features inside the box and the actual object, making classification and localization less reliable \cite{ding2021object,ding2018learning}.

This problem becomes more severe in dense scenes. In harbors, parking lots, airports, or traffic areas, multiple objects are often located close to each other. A single horizontal box may therefore cover not only the target object, but also neighboring instances. As a result, the extracted region features become ambiguous, because they may contain several objects at once. This also affects post-processing: when several horizontal boxes overlap strongly, non-maximum suppression (NMS) can incorrectly remove valid detections of nearby objects \cite{ding2018learning}.

The geometric disadvantage of HBBs can also be illustrated quantitatively. If an oriented rectangle with width $w$, height $h$, and rotation angle $\theta$ is enclosed by its axis-aligned bounding box, the enclosing box has the approximate dimensions

\begin{equation}
    W = w|\cos\theta| + h|\sin\theta|,
    \qquad
    H = w|\sin\theta| + h|\cos\theta|.
\end{equation}

The ratio between the oriented object area and the enclosing horizontal box area is therefore

\begin{equation}
    \mathrm{IoU}_{OBB,HBB}
    =
    \frac{w h}{W H}.
\end{equation}

For an elongated object with an aspect ratio of $5:1$ rotated by $45^\circ$, this ratio drops to approximately $0.28$. This means that more than two thirds of the horizontal box area consists of background. For even more elongated objects, such as ships, the mismatch becomes larger. Therefore, oriented bounding boxes (OBBs) are better suited for aerial imagery, because they can align with the object direction and describe the actual object extent more precisely \cite{ding2021object,ding2018learning}.

\subsubsection{Oriented Bounding Box Parametrization}

An oriented bounding box (OBB) extends the conventional horizontal bounding box by adding an orientation parameter. A common five-parameter representation is defined as

\begin{equation}
    B_{obb} = (x, y, w, h, \theta),
\end{equation}

where $(x,y)$ denotes the center point of the box, $w$ and $h$ describe its width and height, and $\theta$ represents the rotation angle of the rectangle. In contrast to a horizontal bounding box, an OBB can rotate with the object and therefore provides a tighter localization for objects that are not aligned with the image axes \cite{ding2021object,yang2022arbitrary}.

The angle $\theta$ defines the orientation of the bounding box relative to the horizontal image axis. Different works and implementations use different angle ranges. Common choices are $\theta \in [0^\circ, 180^\circ)$ or $\theta \in [-90^\circ, 90^\circ)$. These ranges are sufficient because a rectangle rotated by $180^\circ$ describes the same spatial region as the original rectangle. Therefore, the representation must avoid redundant angle definitions while still describing all possible object orientations \cite{yang2022arbitrary,xie2021oriented}.

A practical difficulty is that different OBB conventions define $w$, $h$, and $\theta$ differently. In an OpenCV-style representation, the angle is tied to a specific rectangle side and the width and height may be swapped depending on the rotation. In contrast, long-edge definitions usually define the angle with respect to the longer side of the rectangle. As a result, the same physical box can have different numerical representations depending on the chosen convention. This is important because inconsistent angle conventions can lead to discontinuities during training and evaluation \cite{yang2022arbitrary}.

For aerial imagery, this parametrization is more precise than horizontal bounding boxes because many objects, such as ships, airplanes, vehicles, bridges, and storage tanks, appear with arbitrary orientations. By aligning the bounding box with the object direction, OBBs reduce the amount of included background and better separate neighboring objects in dense scenes. This makes them especially useful for oriented object detection in aerial and satellite images \cite{ding2021object,ding2018learning,xie2021oriented}.

\subsubsection{Angle Regression Challenges}

Although oriented bounding boxes provide a more precise object representation for aerial imagery, their angle parameter introduces additional regression challenges. A common problem is the boundary discontinuity of angle regression. Since rectangular boxes are periodic, an orientation of $0^\circ$ and $180^\circ$ can describe the same spatial direction, but a standard regression loss treats them as numerically far apart. For example, a prediction of $179^\circ$ for a ground-truth angle of $0^\circ$ may be geometrically almost correct, while an L1 or Smooth-L1 loss produces a large error. This causes loss spikes near the angle boundary and makes training unstable \cite{yang2021learning}.

A second issue is the ambiguity of width, height, and angle. For square-like objects, a rotation by $90^\circ$ can result in nearly the same bounding box, because width and height are almost interchangeable. Even for rectangular objects, different angle conventions may swap $w$ and $h$ at certain boundary cases. This means that two numerically different parameter sets can describe almost identical boxes. As a result, directly regressing $(x, y, w, h, \theta)$ can lead to discontinuous targets and inconsistent gradients during training \cite{yang2021learning}.

One solution is Circular Smooth Label (CSL), which reformulates angle prediction as a classification task instead of a direct regression task. In this approach, the continuous angle space is divided into discrete bins, and the periodic nature of angles is modeled by circular labels. Therefore, angles close to each other across the boundary, such as $0^\circ$ and $179^\circ$, are also treated as neighboring classes. This reduces the boundary discontinuity problem and increases the tolerance to small angular deviations \cite{yang2021learning}.

Another line of work avoids direct angle regression by representing rotated bounding boxes as Gaussian distributions. Instead of comparing the five box parameters independently with an L1-type loss, the predicted and ground-truth boxes are converted into two-dimensional Gaussian distributions. The distance between both boxes can then be measured using Kullback-Leibler Divergence (KLD). This couples position, size, and orientation in a single loss function and can adapt the importance of angle errors depending on the object shape. This is especially useful for elongated objects such as ships or vehicles, where small angle errors can strongly reduce the localization quality \cite{yang2021learning}.

Overall, angle prediction is difficult because the numerical representation of an oriented box is not always continuous, unique, or equally sensitive for all object shapes. CSL addresses this problem by treating angle prediction as circular classification, while KLD-based Gaussian representations address it by comparing whole boxes as probability distributions rather than isolated parameters. Both approaches aim to make oriented object detection more stable and better aligned with the geometric properties of rotated objects \cite{yang2021learning}.


\subsubsection{Key Models}

One important model for oriented object detection in aerial imagery is the RoI Transformer. It is based on a two-stage detection framework and addresses the mismatch between horizontal region proposals and arbitrarily oriented objects. Instead of directly relying on horizontal Regions of Interest (HRoIs), the RoI Transformer learns a transformation from HRoIs to rotated Regions of Interest (RRoIs). This is done by a trainable RRoI Learner, which predicts the geometry of a rotated proposal from the features of a horizontal proposal. The resulting RRoIs are then used for rotation-aware feature extraction through Rotated Position-Sensitive RoI Align. This improves the alignment between region features and object geometry, especially in dense aerial scenes where horizontal boxes may contain several neighboring objects \cite{ding2018learning}.

Oriented R-CNN is another representative two-stage model for oriented object detection. Its main goal is to generate high-quality oriented proposals more efficiently than previous proposal-based approaches. While rotated RPNs use many rotated anchors with different scales, aspect ratios, and angles, and RoI Transformer learns oriented proposals through an additional transformation step, Oriented R-CNN introduces a lightweight oriented RPN. This oriented RPN predicts oriented proposals directly by using a midpoint offset representation. The second stage then applies an oriented R-CNN head with rotated RoIAlign to classify the proposals and refine their spatial locations \cite{xie2021oriented}.

Compared with RoI Transformer, Oriented R-CNN focuses more strongly on reducing the computational cost of proposal generation. Xie et al. argue that proposal generation is a major bottleneck in many two-stage oriented detectors. By using a simpler oriented RPN, Oriented R-CNN can generate oriented proposals in a nearly cost-free manner while maintaining strong detection accuracy on oriented object detection benchmarks such as DOTA and HRSC2016 \cite{xie2021oriented}.

Besides proposal-based methods, one-stage and anchor-free oriented detectors represent another important direction. These methods aim to predict oriented bounding boxes more directly, often without a separate proposal generation stage. In principle, this can reduce the complexity of anchor design and avoid the need to define many rotated anchors with different angles. However, with the literature used in this work, the main focus remains on proposal-based oriented detectors, especially RoI Transformer and Oriented R-CNN, because they are directly covered by the selected sources \cite{ding2018learning,xie2021oriented}.


\subsection{Datasets and Evaluation}

\subsubsection{Benchmark Datasets}
% - DOTA (Dataset for Object deTection in Aerial images):
%   * Largest and most widely used aerial detection benchmark
%   * OBB annotations, 15 object categories (vehicles, ships, planes, ...)
%   * DOTA-v1.0, v1.5, v2.0 with increasing size and categories
% - DIOR (Object Detection in Optical Remote sensing):
%   * 20 categories, balanced distribution, simpler annotations (HBB)
%   * Good for baseline comparisons
% - xView:
%   * Real satellite imagery, very small objects, 60 categories
%   * Challenging due to extreme scale and resolution variation

\subsubsection{Evaluation Metrics}
% - mAP (mean Average Precision) recap: area under PR curve, averaged over classes
% - IoU for OBB differs from standard axis-aligned IoU:
%   * Rotated rectangle intersection is non-trivial to compute
%   * Approximations vs. exact polygon intersection
% - Typical threshold: mAP@0.5 (IoU > 0.5 counts as true positive)
% - Some works report mAP@0.5:0.95 for stricter evaluation


%%%-----------------------------------------------------------------------------

\section{Exercises and Homeworks}



%%%-----------------------------------------------------------------------------

\section*{References}

\printbibliography[heading=noheader]

%%%-----------------------------------------------------------------------------
\end{document}
%%%-----------------------------------------------------------------------------
